\documentclass[reqno, 12pt]{amsart}

\usepackage{amssymb}
\usepackage{amsthm}
\usepackage{amsmath}
\usepackage{amsxtra}
\usepackage{mathrsfs}
\usepackage{comment}
\usepackage{graphicx}
\usepackage{placeins}
\usepackage{multicol}
\usepackage{bbm}
\usepackage{stmaryrd}
%\usepackage{enumerate}
\usepackage[inline, shortlabels]{enumitem}
\usepackage{mathdots}
\usepackage{colonequals}
\usepackage{hyperref}


\usepackage{calc}
\usepackage{tikz}
\usetikzlibrary{arrows,calc,automata,shadows,backgrounds,positioning,intersections,fadings,decorations.pathreplacing,shapes,matrix}
\usepackage{tikz-cd}

\usepackage{fullpage}

\newtheorem{thm}{Theorem}[section]
\newtheorem{lem}[thm]{Lemma}
\newtheorem{cor}[thm]{Corollary}
\newtheorem{conj}[thm]{Conjecture}
\newtheorem{prop}[thm]{Proposition}
\theoremstyle{definition}  
\newtheorem{defn} [thm] {Definition} 
\newtheorem{claim}[thm]{Claim}
\newtheorem{example} [thm] {Example}
\newtheorem{rem} [thm] {Remark}

\newenvironment{problab}[1]
{\noindent\textbf{Problem #1}.}
{\vskip 6pt}

\newenvironment{exolab}[1]
{\noindent\textbf{Exercise #1}.}
{\vskip 6pt}

\theoremstyle{remark}
\newtheorem*{soln}{Solution}

\newcommand{\C}{\mathbb C}
\renewcommand{\H}{\mathbb H}
\newcommand{\D}{\mathbb D}
\newcommand{\F}{\mathbb F}
\newcommand{\Q}{\mathbb Q}
\newcommand{\R}{\mathbb R}
\newcommand{\Z}{\mathbb Z}
\newcommand{\T}{\mathbb T}
\renewcommand{\P}{\mathcal P}
\renewcommand{\O}{\mathcal O}
\newcommand{\m}{\mathfrak m}
\newcommand{\A}{\mathbb A}
\renewcommand{\SS}{\mathbb S}
\renewcommand{\L}{\mathcal L}

\newcommand{\B}{\mathcal B}
\newcommand{\E}{\mathcal E}

\newcommand{\SSpan}{\text{span}}
\newcommand{\la}{\langle}
\newcommand{\ra}{\rangle}

\DeclareMathOperator{\lcm}{lcm}
\DeclareMathOperator{\img}{img}
\DeclareMathOperator{\Ann}{Ann}
\DeclareMathOperator{\Tor}{Tor}
\DeclareMathOperator{\tr}{tr}
\DeclareMathOperator{\Hom}{Hom}
\DeclareMathOperator{\End}{End}
\DeclareMathOperator{\Aut}{Aut}
\DeclareMathOperator{\Gal}{Gal}
\DeclareMathOperator{\sgn}{sgn}
\DeclareMathOperator{\ord}{ord}
\DeclareMathOperator{\ad}{ad}
\DeclareMathOperator{\coker}{coker}
\DeclareMathOperator{\rank}{rank}
\DeclareMathOperator{\minpoly}{minpoly}

\DeclareMathOperator{\id}{id}

\usepackage{mathpazo}

\everymath{\displaystyle}

\tikzset{commutative diagrams/.cd, arrow style = tikz, diagrams = {>=latex}}

\newenvironment{amatrix}[1]{%
  \left(\begin{array}{@{}*{#1}{c}|c@{}}
}{%
  \end{array}\right)
}

\usepackage{abstract}
\renewcommand{\abstractname}{}    % clear the title
\renewcommand{\absnamepos}{empty} % originally center
%\setlength{\abstitleskip}

\begin{document}
%\renewcommand{\abstractname}{\vspace{-\baselineskip}}

\title{18.700 Problem Set 8}
%\author{Évariste Galois}
\dedicatory{Due Wednesday, November 27 at 11:59 pm on \href{https://canvas.mit.edu/courses/27315}{Canvas}}
%\thanks{}
%\contrib{}
%\address{}

\maketitle

\begin{abstract}
  \noindent Collaborated with:\\
  Sources used: 
\end{abstract}
\vskip 0.5cm

Let $V$ and $W$ be nonzero finite-dimensional inner product spaces over $\F$.
\vskip 0.5cm

\begin{problab}{1} (7A \#4) (7 points)
  Suppose $T \in \L(V)$ and $U$ is a subspace of $V$. Prove that
  \begin{center}
    $U$ is $T$-invariant $\iff$ $U^\perp$ is $T^*$-invariant.
  \end{center}
\end{problab}

%\begin{soln}
%\end{soln}

\begin{problab}{2} (7A \#5) (6 points)
  Suppose $T \in \L(V,W)$. Suppose $e_1, \ldots, e_n$ is an orthonormal basis of $V$ and $f_1, \ldots, f_m$ is an orthonormal basis of $W$. Prove that
  \begin{equation*}
    \|T(e_1)\|^2 + \cdots + \|T(e_n)\|^2 = \|T^*(f_1)\|^2 + \cdots + \|T^*(f_m)\|^2 \, .
  \end{equation*}
\end{problab}

\begin{problab}{3} (7B \#1) (7 points)
  Let $\F = \C$ and suppose $T \in \L(V)$ is normal. Show that $T$ is self-adjoint if and only if all the eigenvalues of $T$ are real.
\end{problab}

\begin{problab}{4} (7B \#6) (6 points)
  Let $\F = \C$ and suppose $T \in \L(V)$ is a normal operator such that $T^9 = T^8$. Prove that $T$ is self-adjoint and $T^2 = T$. (\textit{Hint}: Use the previous exercise.)
\end{problab}

\begin{problab}{5} (7C \#3) (6 points)
  Let $n$ be a positive integer and $T \in \L(\F^n)$ be the operator whose matrix with respect to the standard basis consists of all $1$s. Show that $T$ is a positive operator.
\end{problab}

\begin{problab}{6} (7E \#3) (4 points)
  Give an example of $T \in \L(\C^2)$ such that $0$ is the only eigenvalue of $T$ and the singular values of $T$ are $5,0$.
\end{problab}

\end{document} 	
